% 02-introduction.tex
% Section 2 — Introduction. Drafted G3.6. [ERIC-REVIEW] on positioning,
% state-of-the-art references, and contribution framing.

\section{Introduction}
\label{sec:introduction}

\subsection{The Collatz conjecture}

The Collatz conjecture, attributed to Lothar Collatz around 1937,
concerns the behaviour of the map
\[
\collatz : \Nat \to \Nat, \qquad
\collatz(n) = \begin{cases} n / 2 & \text{if } n \text{ is even,}\\ 3n + 1 & \text{if } n \text{ is odd.}\end{cases}
\]
The conjecture asserts that, for every positive integer \(n\), the orbit
\(n, \collatz(n), \collatz^2(n), \ldots\) eventually reaches \(1\).
Equivalently, one can show that the conjecture is equivalent to the
simultaneous absence of two phenomena : divergent orbits, and
\emph{nontrivial cycles} (periodic orbits different from the loop
\(1 \mapsto 4 \mapsto 2 \mapsto 1\)).

% [ERIC-REVIEW] Historical attribution "around 1937" — please confirm or
% refine. Some sources attribute earlier/later. Non-critical but should
% be accurate.

\subsection{State of the art on cycle exclusion}

The \emph{cycle problem} has attracted sustained attention. Representative
recent milestones :

\begin{itemize}
  \item Barina~\cite{Barina2025} has verified computationally that every
        positive integer below \(2^{71}\) reaches \(1\) under
        \(\collatz\). This in particular excludes any nontrivial cycle
        whose minimum value is below that threshold.
  \item Hercher~\cite{Hercher2023} has shown that no Collatz \(m\)-cycle
        exists for \(m \leq 91\), building on prior
        Simons--de Weger--type analysis~\cite{SimonsDeWeger2005}.
  \item Several formalisations in interactive theorem provers have
        appeared recently, targeting partial results on the cycle
        problem with varying scope and rigour levels. The present work
        differs from these efforts in its specific conditional form and
        in the explicit management of its axiom baseline
        (see Section~\ref{sec:lean}).
\end{itemize}

% [ERIC-REVIEW] The "several formalisations" sentence is deliberately
% vague ; if you want to name specific prior formalisations by other
% authors (cf. your memory of e.g. johnjanik, fc0web, kaiarasanen) please
% provide the names and we'll cite them factually. Current phrasing is
% safe but less informative.

Against this backdrop, the present paper does \emph{not} claim to resolve
the Collatz conjecture nor the cycle problem unconditionally. It
formalises, in Lean~4, a conditional non-existence of nontrivial odd
cycles, under three explicit hypotheses (Section~\ref{sec:hypotheses}).

\subsection{Contributions}

% [ERIC-REVIEW] contributions phrasing. Please confirm these are
% accurate and appropriately modest.

The main contributions of this paper and the accompanying repository are
the following :

\begin{enumerate}
  \item A Lean~4 formalisation of a conditional cycle-exclusion theorem
        (Theorem~\ref{thm:main}) whose central axiom chain is exactly
        \([\texttt{propext},\, \texttt{Classical.choice},\,
        \texttt{Quot.sound}]\), the three standard kernel axioms of
        Mathlib.
  \item An explicit, auditable separation between \emph{published}
        hypotheses (\textsf{BakerSeparation} from~\cite{Baker1966}, and
        \textsf{BarinaVerification} from~\cite{Barina2025}) and a
        \emph{derived} hypothesis (\textsf{DerivedLargeKBound}),
        formalised as named Lean structures rather than as
        user-declared axioms.
  \item An end-to-end reproducibility pipeline (\texttt{reproduce.sh}
        with five exit codes, matching probes in
        \texttt{probes/check\_central\_axioms.lean} and
        \texttt{probes/check\_sorry.lean}, and a continuous integration
        workflow) that enforces the axiom baseline on every commit.
  \item A public audit trail under \texttt{docs/BIBLE/} documenting the
        gating decisions, risk register, limitations, and adversarial
        reviews that produced the current state of the repository.
\end{enumerate}

\subsection{Organisation}

The remainder of the paper is organised as follows.
Section~\ref{sec:preliminaries} fixes notation and defines the formal
predicate \texttt{IsOddCycle} used throughout.
Section~\ref{sec:main} states the main theorem
and sketches its proof structure. Section~\ref{sec:hypotheses} details
the three hypotheses, distinguishing published from derived components.
Section~\ref{sec:lean} describes the Lean~4 formalisation, including the
axiom baseline. Section~\ref{sec:reproducibility} gives the
reproducibility pipeline. Section~\ref{sec:conclusion} identifies the
main item of future work (Phase Legendre : promoting
\textsf{DerivedLargeKBound} from a structure to a proven theorem) and
discusses implications.

% [ERIC-REVIEW] If desired, we can add a brief "acknowledgements" stub
% here for potential reviewers / colleagues. Currently omitted.
